\section{Desenvolvimento e implementação}
\label{sec:desenvolvimento}

\subsection{Base de execução}

O projeto utiliza Python 3.12 e empacotamento definido em \code{pyproject.toml}. FastAPI fornece os endpoints e o ciclo de vida ASGI; Uvicorn executa as aplicações; HTTPX realiza as chamadas entre processos; Pydantic valida os contratos; Pytest e pytest-asyncio compõem a suíte de testes \cite{fastapi,uvicorn,httpx,pydantic,pytest,pytestasyncio}.

O \code{Dockerfile} \cite{dockerfile} produz uma imagem comum. No Compose, \code{node1} realiza o build e os demais serviços reutilizam a imagem. Cada nó recebe uma lista de peers com os endereços internos \code{http://nodeX:8000}. As portas 8001, 8002, 8003 e 8010 são apenas mapeamentos para acesso pelo host.

\subsection{Estado local e sincronização entre corrotinas}

A classe \code{NodeState} concentra o estado mutável e usa um único \code{asyncio.Lock}. A documentação de Python define \code{Lock} como uma primitiva de exclusão mútua entre tarefas assíncronas e \code{Event} como uma sinalização para tarefas em espera \cite{python-asyncio}. No projeto, o lock protege:

\begin{itemize}
    \item valor do relógio e contadores de mensagens;
    \item estado, timestamp e respostas do mutex;
    \item líder conhecido, rodada e respostas de eleição;
    \item histórico circular de até 500 registros.
\end{itemize}

Um \code{asyncio.Event} é criado em cada pedido de mutex. Ele é sinalizado quando o conjunto de respostas atinge o número de peers. O endpoint usa \code{asyncio.wait_for()} para limitar a espera, sem implementar polling ativo \cite{python-waitfor}.

\subsection{Relógio de Lamport e log de eventos}

\subsubsection{Operações do relógio}

O módulo \code{clock.py} mantém um contador inteiro iniciado em zero. As operações são pequenas e determinísticas:

\begin{lstlisting}[language=Python,caption={Operações essenciais do relógio lógico.}]
def tick(self) -> int:
    self.value += 1
    return self.value

def update(self, received_time: int) -> int:
    self.value = max(self.value, received_time) + 1
    return self.value
\end{lstlisting}

Um envio executa \code{tick()}, cria a mensagem com o valor obtido e registra \code{SEND} com o mesmo timestamp. Um recebimento executa \code{update()} antes de registrar \code{RECEIVE}. Assim, o timestamp local do recebimento é sempre maior que o valor transportado.

\subsubsection{Eventos e anotações}

O histórico não é uma cópia literal da definição matemática de evento. Para separar a linha causal principal dos detalhes do protocolo, o modelo distingue:

\begin{itemize}
    \item \code{lamport_event}: evento central usado na linha causal, como \code{LOCAL}, \code{SEND}, \code{RECEIVE}, \code{MUTEX_WANTED}, \code{ENTER_CS}, \code{EXIT_CS}, \code{ELECTION_STARTED} e \code{LEADER_TIMEOUT};
    \item \code{annotation}: detalhe ou transição auxiliar, como uma resposta adiada, uma mudança de líder, a rejeição de coordenador ou uma falha de envio.
\end{itemize}

Por exemplo, o recebimento de um pedido de mutex gera um \code{RECEIVE} e uma anotação \code{MUTEX_REQUEST_RECEIVED} com o mesmo timestamp. Isso representa um único evento de recebimento e duas visões de observabilidade, não dois eventos sucessivos do processo. A classificação é semântica, não um mecanismo de controle do relógio: a transição \code{MUTEX_HELD}, por exemplo, executa \code{tick()} e ainda assim é rotulada como \code{annotation}. Portanto, reconstruções precisas devem considerar a ação e o timestamp, e não apenas \code{event_kind}.

\subsubsection{Exemplo causal}

Considere o encadeamento usado no smoke test:

\begin{enumerate}
    \item \code{node1} registra um evento local;
    \item \code{node1} envia uma mensagem a \code{node2};
    \item \code{node2} recebe a mensagem e depois envia outra a \code{node3};
    \item \code{node3} recebe a segunda mensagem.
\end{enumerate}

O script verifica as desigualdades:

\begin{align*}
    C(S_{1\rightarrow2}) &< C(R_{1\rightarrow2}),\\
    C(R_{1\rightarrow2}) &< C(S_{2\rightarrow3}),\\
    C(S_{2\rightarrow3}) &< C(R_{2\rightarrow3}).
\end{align*}

Essas relações demonstram a condição de Lamport no cenário observado. Elas não provam que qualquer par de timestamps ordenados possua relação causal.

\subsection{Exclusão mútua distribuída}

\subsubsection{Modelagem de estados}

A implementação representa o estado local por três valores:

\begin{table}[H]
    \centering
    \small
    \begin{tabularx}{\textwidth}{@{}p{2.8cm}X X@{}}
        \toprule
        Estado & Significado & Comportamento ao receber pedido \\
        \midrule
        \code{RELEASED} & Não deseja e não ocupa a seção crítica. & Responde imediatamente. \\
        \code{WANTED} & Possui pedido local aguardando respostas. & Compara prioridades e pode adiar. \\
        \code{HELD} & Está dentro da seção crítica. & Sempre adia a resposta. \\
        \bottomrule
    \end{tabularx}
    \caption{Estados locais usados para Ricart--Agrawala.}
    \label{tab:mutex-states}
\end{table}

Esses nomes são a modelagem adotada no projeto; a referência acadêmica é usada para as regras do algoritmo, não para afirmar que a nomenclatura seja textual no artigo original.

\subsubsection{Criação e disseminação do pedido}

Ao iniciar uma tentativa, o nó:

\begin{enumerate}
    \item exige estado \code{RELEASED};
    \item incrementa o relógio e fixa esse valor em \code{request_timestamp};
    \item gera um \code{request_id} único;
    \item muda para \code{WANTED};
    \item limpa respostas antigas;
    \item envia \code{MUTEX_REQUEST} a todos os peers.
\end{enumerate}

O \code{request_timestamp} é igual para todos os pedidos da mesma tentativa. Já o \code{logical_time} do envelope muda a cada envio, pois cada transmissão é um evento distinto. Essa separação é necessária: a prioridade pertence à tentativa, e não ao instante em que cada peer foi contatado.

\subsubsection{Regra de prioridade}

A função pura de prioridade usa comparação lexicográfica:

\begin{equation}
    local\_priority \Leftrightarrow
    (t_{local}, id_{local}) < (t_{remoto}, id_{remoto}).
\end{equation}

Se o estado local for \code{WANTED} e o pedido local tiver prioridade, a resposta remota é armazenada em \code{deferred_replies}. Caso contrário, \code{MUTEX_REPLY} é enviada imediatamente. Conjuntos impedem que respostas duplicadas sejam contadas mais de uma vez, e o \code{request_id} impede que uma resposta antiga libere uma tentativa nova.

\subsubsection{Entrada, saída e abort}

A passagem para \code{HELD} só ocorre depois de receber exatamente uma resposta de cada peer. A entrada e a saída da seção crítica são eventos locais separados. Após a saída, o estado volta a \code{RELEASED} e todas as respostas adiadas são enviadas.

Se a espera ultrapassa o timeout, a tentativa é encerrada por \code{MUTEX_ABORTED}. Essa ação incrementa o relógio, limpa o estado local e devolve as respostas adiadas. O procedimento melhora a controlabilidade dos testes, porém não garante que todos os processos tenham visão consistente do abort nem resolve falha de participante.

\subsubsection{Observador externo}

Antes e depois do período simulado na seção crítica, o nó envia notificações a \code{resource}. O observador mantém um mapa de solicitações ativas. Uma nova entrada enquanto o mapa não está vazio incrementa \code{violations}. Esse mecanismo fornece evidência de ausência de sobreposição nos cenários executados, mas não faz parte da autorização e não é uma prova formal para todas as execuções possíveis.

\subsection{Eleição de líder}

\subsubsection{Mensagens e fluxo}

A Tabela~\ref{tab:election-messages} relaciona as mensagens com sua função.

\begin{table}[H]
    \centering
    \small
    \begin{tabularx}{\textwidth}{@{}p{3.5cm}X p{3.7cm}@{}}
        \toprule
        Tipo & Função & Origem \\
        \midrule
        \code{ELECTION} & Pergunta a processos de ID maior se algum está ativo. & Regra central do Bully. \\
        \code{ELECTION_OK} & Confirma que um processo superior está ativo. & Nome local para resposta superior. \\
        \code{COORDINATOR} & Anuncia o vencedor da eleição. & Regra central do Bully. \\
        \code{HEARTBEAT} & Renova a presença do líder conhecido. & Adaptação do projeto. \\
        \bottomrule
    \end{tabularx}
    \caption{Mensagens do subsistema de eleição.}
    \label{tab:election-messages}
\end{table}

Ao iniciar uma rodada, o nó gera um UUID em \code{election_id}, incrementa o relógio e seleciona apenas peers de ID maior. Se a lista estiver vazia, torna-se líder e anuncia o resultado. Se nenhum superior responder dentro de 700 ms, o iniciador também se declara líder. Quando recebe ao menos uma resposta, espera até 2500 ms por \code{COORDINATOR}; na ausência do anúncio, registra \code{ELECTION_ROUND_TIMED_OUT} e inicia uma nova rodada.

\subsubsection{Detecção de falha e recuperação}

O líder envia heartbeat a cada 700 ms. Seguidores usam \code{time.monotonic()} para medir o tempo desde a última mensagem aceita. Após 2500 ms, removem o líder, registram \code{LEADER_TIMEOUT} e agendam uma eleição. No startup, cada nó também agenda uma eleição após 500 ms; isso permite que um processo superior recuperado volte a disputar a liderança.

\subsubsection{Aceitação de líder}

A aceitação de um coordenador não depende apenas da existência de uma eleição em andamento; ela respeita a prioridade dos identificadores. A regra implementada é:

\begin{equation}
\begin{split}
    accept(announced) \Leftrightarrow {} & announced \geq node\_id \\
    & \land (current = \varnothing \lor announced \geq current).
\end{split}
\label{eq:accept-leader}
\end{equation}

A primeira condição impede que um processo ativo aceite como líder alguém inferior a si próprio. A segunda impede rebaixar um líder já conhecido. O heartbeat exige ainda que \code{sender_id == leader_id}.

Se um coordenador ou heartbeat coerente, porém inferior ao nó local, é rejeitado, a camada de controle agenda uma eleição. O agendamento usa uma task única protegida por lock, evitando a criação simultânea de várias tarefas locais. Mensagens malformadas, como \code{leader_id} diferente do remetente, são rejeitadas sem serem tratadas como evidência de uma liderança inferior válida.

\subsubsection{Limites do identificador de eleição}

O \code{election_id} é associado à rodada local do iniciador e permite ignorar \code{ELECTION_OK} duplicado, atrasado ou referente a outra tentativa. Ele não é um número de época global nem identifica todas as eleições concorrentes do sistema. A convergência depende da regra de prioridade, do anúncio do coordenador e dos timeouts.

\subsection{Decisões de implementação e limitações}

As decisões centrais foram manter o estado em memória, evitar qualquer coordenador para o mutex, reutilizar um envelope comum e realizar a falha do líder pela parada real do contêiner. Como consequência:

\begin{itemize}
    \item reiniciar um nó perde seu histórico e estado local;
    \item falhas durante uma rodada de mutex podem interromper o progresso;
    \item os timeouts foram calibrados para a execução local;
    \item uma partição pode criar visões incompatíveis de liderança;
    \item o serviço \code{resource} pode detectar sobreposição observada, mas não impedir uma violação;
    \item a criação de um novo \code{httpx.AsyncClient} por chamada é simples, porém não aproveita pooling de longo prazo sugerido pela documentação do HTTPX \cite{httpx}.
\end{itemize}
